\chapter{Central series and finite nilpotent groups}\label{chap:M07.18}
The characteristic-subgroup definition, commutator convention and solvability
theorems remain in \ref{K07.04.characteristic_subgroups},
\ref{K07.04.commutator_abelianization} and \ref{K07.15.solvable_closure}.
Here $[x,y]=xyx^{-1}y^{-1}$. Groups are arbitrary until finiteness is stated;
$p$ always denotes a prime. No classification of finite simple groups is used.

\KEntry{K07.18.mixed_commutators}{Commutators of subgroups}
For subgroups $A,B\leq G$, let $[A,B]$ be the subgroup generated by all
$[a,b]$ with $a\in A,b\in B$. Then $[A,B]=[B,A]$, since
$[a,b]^{-1}=[b,a]$. Inclusion in either argument preserves inclusion of
these generated groups. For a homomorphism $f$,
\[
 f([A,B])=[f(A),f(B)].
\]
Indeed $f([a,b])=[f(a),f(b)]$, and the image-of-generated-subgroup
theorem \ref{K07.01.hom_on_generators} proves equality, including its
surjectivity onto the displayed generated image. For an automorphism
$\alpha$, in particular, $\alpha([A,B])=[\alpha(A),\alpha(B)]$.
Thus $[A,B]$ is characteristic in $G$ if both $A,B$ are characteristic,
and is normal in $G$ if both are normal: apply this identity to all
automorphisms, or just to inner automorphisms, respectively.
If $A$ is normal, $[a,g]=a(ga^{-1}g^{-1})\in A$, so
$[A,G]\leq A$. Finally, for $N\trianglelefteq G$,
\[
 [G,A]\leq N\quad\Longleftrightarrow\quad
 AN/N\leq Z(G/N).
\]
Both sides say that every image of $a\in A$ commutes with every image
of $g\in G$. This criterion also includes $A=1$ or $N=G$.

\KEntry{K07.18.central_series}{Upper and lower central series}
Define
\[
 Z_0(G)=1,\qquad Z_{i+1}(G)/Z_i(G)=Z(G/Z_i(G)),
 \qquad \gamma_1(G)=G,\quad\gamma_{i+1}(G)=[\gamma_i(G),G].
\]
The first definition means full inverse image under the quotient map.
Each $Z_i$ and $\gamma_i$ is characteristic in $G$; the upper series
increases and the lower series decreases. To prove this without assuming
the quotients exist in advance, start with characteristic $Z_0$.
An automorphism preserving $Z_i$ induces an automorphism of $G/Z_i$,
which preserves its center, because commuting with every element is
unchanged under an automorphism. Its inverse gives equality, so the
center's full inverse image $Z_{i+1}$ is characteristic and hence normal.
The inverse image contains $Z_i$, proving increase. For the lower series,
the preceding commutator identity proves characteristicity inductively;
normality and $[\gamma_i,G]\leq\gamma_i$ prove decrease. Thus also
\[
 [G,Z_{i+1}]\leq Z_i.
\]
If two consecutive upper terms agree, the corresponding quotient has
trivial center, and the definition makes every later term equal to them.
If two consecutive lower terms agree, applying $[-,G]$ repeatedly gives
the same conclusion for that series. In a finite group both chains
eventually stabilize by strict changes of subgroup orders. Stabilization
alone does not say that their endpoints are $G$ and $1$.

\KEntry{K07.18.nilpotence_criteria}{Nilpotence and central-series comparison}
A group is nilpotent of class at most $c$, where $c\geq0$, if
$Z_c(G)=G$. Its class is the smallest such $c$. The trivial group has
class zero, and a nontrivial abelian group has class one.
\begin{theorem}\label{thm:central-series-criteria}
The following are equivalent: $Z_c(G)=G$; $\gamma_{c+1}(G)=1$; and
there is a chain of normal subgroups
\[
 1=C_0\leq C_1\leq\cdots\leq C_c=G,
 \qquad [G,C_i]\leq C_{i-1}\quad(1\leq i\leq c).
\]
For any such chain, $C_i\leq Z_i(G)$ and
$\gamma_j(G)\leq C_{c+1-j}$ for $1\leq j\leq c+1$.
In particular, if $G$ has class at most $c$, then
$\gamma_{c+1-i}(G)\leq Z_i(G)$ for $0\leq i\leq c$.
\end{theorem}
\begin{proof}
For a central chain, induct on $i$. The case $C_0\leq Z_0$ is equality.
If $C_{i-1}\leq Z_{i-1}$, then
$[G,C_i]\leq Z_{i-1}$, so the image of $C_i$ in $G/Z_{i-1}$ is
central and $C_i\leq Z_i$. For the lower comparison start with
$\gamma_1=G=C_c$; inclusion in the commutator arguments gives
$\gamma_{j+1}\leq[C_{c+1-j},G]\leq C_{c-j}$.
The upper series itself is a central chain if $Z_c=G$, and therefore
forces $\gamma_{c+1}=1$. Conversely, if $\gamma_{c+1}=1$, set
$C_i=\gamma_{c+1-i}$. The defining commutator recurrence makes this
a central chain, and its upper comparison forces $G=C_c\leq Z_c$.
Apply the lower comparison to the upper chain for the last inclusion.
For $c=0$ all three assertions say $G=1$.
\end{proof}
Consequently the least upper and lower termination lengths give the same
class. There is no general reverse inclusion interlacing all upper and
lower terms; an explicit counterexample appears below.

\KEntry{K07.18.nilpotence_closure}{Subgroups, images and finite products}
\begin{proposition}\label{prop:nilpotence-closure}
A subgroup or homomorphic image of a nilpotent group of class at most
$c$ is nilpotent of class at most $c$. A finite product of nilpotent
groups is nilpotent, of class the maximum of the factor classes; the empty
product has class zero.
\end{proposition}
\begin{proof}
For $H\leq G$, induction and commutator monotonicity give
$\gamma_i(H)\leq\gamma_i(G)$. For a surjection $f:G\to Q$ the
commutator image identity gives $f(\gamma_i(G))=\gamma_i(Q)$ by
induction, starting with $f(G)=Q$. Apply the termination criterion.
In $A\times B$, commutators are coordinatewise. The commutators of
$(a,1),(a',1)$ and $(1,b),(1,b')$ give the two coordinate generated
subgroups separately, so induction gives
$\gamma_i(A\times B)=\gamma_i(A)\times\gamma_i(B)$.
Termination therefore has length the maximum of the two lengths.
Induct over a finite list of factors. The zero-factor convention is the
one-element group. No assertion is made here about an infinite product
whose factor classes have no common finite bound.
\end{proof}

\KEntry{K07.18.pgroup_nilpotent}{Finite prime-power groups are nilpotent}
\begin{theorem}\label{thm:pgroup-nilpotent}
If $|P|=p^a$ with $a\geq0$, then $P$ has class at most $a$.
For $a\geq2$ the sharper bound $a-1$ holds. For $a=1$ its class is
one, and for $a=0$ its class is zero.
\end{theorem}
\begin{proof}
Each nontrivial quotient $P/Z_i(P)$ is a finite $p$-group and has
nontrivial center by \ref{K07.06.pgroup_center}. Its center has at least
$p$ elements by Lagrange. Thus every strict upper-series step multiplies
the order by at least $p$, and $a$ such steps exhaust $P$.
If $a\geq2$ and the class were exactly $a$, all these steps would have
index exactly $p$. The quotient $P/Z_{a-2}$ would have order $p^2$,
so would be abelian by \ref{K07.06.prime_square_classification}.
Its entire group would be its center, forcing $Z_{a-1}=P$, a
contradiction. A group of prime order is nontrivial cyclic, so its class
is one; the order-one case is trivial. These cases explain why the
sharper bound must not be stated without $a\geq2$.
\end{proof}

\KEntry{K07.18.normalizer_condition}{The normalizer condition}
\begin{theorem}\label{thm:nilpotent-normalizer}
If $G$ is nilpotent and $H<G$, then $H<N_G(H)$. Finiteness of $G$
or $H$ is not needed.
\end{theorem}
\begin{proof}
Choose the least $j$ with $Z_j\nsubseteq H$, which exists because
some $Z_c=G$. Then $j\geq1$ and $Z_{j-1}\leq H$.
Choose $z\in Z_j\setminus H$. For every $h\in H$,
$[z,h]\in Z_{j-1}\leq H$, so $zhz^{-1}\in H$.
The same argument with $z^{-1}\in Z_j$ gives the reverse inclusion;
thus $zHz^{-1}=H$. Therefore $z\in N_G(H)\setminus H$.
\end{proof}

\KEntry{K07.18.sylow_normalizer}{Normalizer of a Sylow normalizer}
\begin{lemma}\label{lem:sylow-normalizer-self}
For a Sylow subgroup $P$ of a finite group $G$, put $N=N_G(P)$.
Then $N_G(N)=N$.
\end{lemma}
\begin{proof}
The subgroup $P$ is normal in $N$ by its definition, and it is a Sylow
subgroup of $N$ because its order already has the full $p$-part of
$|G|$. Sylow conjugacy makes it the unique Sylow $p$-subgroup of $N$.
An automorphism of $N$ preserves subgroup orders and therefore fixes
$P$, so $P$ is characteristic in $N$. If $g$ normalizes $N$,
conjugation by $g$ is an automorphism of $N$, hence fixes $P$.
Thus $g\in N_G(P)=N$. Conversely every subgroup normalizes itself,
so $N\leq N_G(N)$, proving equality.
\end{proof}

\KEntry{K07.18.finite_nilpotent}{Finite nilpotence and Sylow products}
\begin{theorem}\label{thm:finite-nilpotent-sylow}
For a finite group $G$ the following are equivalent:
$G$ is nilpotent; every proper subgroup is properly contained in its
normalizer; every Sylow subgroup is normal; and multiplication gives an
isomorphism from the product of one Sylow subgroup for each prime
dividing $|G|$ onto $G$.
\end{theorem}
\begin{proof}
Nilpotence implies the normalizer condition by the preceding theorem.
Assume that condition and take a Sylow $P$. Its normalizer $N$ is
self-normalizing by the lemma. If $N<G$ the condition would force
$N<N_G(N)$, impossible. Hence $N=G$ and $P$ is normal.

Now suppose all Sylow subgroups $P_1,\ldots,P_s$ are normal, with
distinct underlying primes. Inductively the product
$H=P_1\cdots P_{t-1}$ is normal, is the internal direct product of
those factors, and has their product order; start with $H=1$.
Its order is coprime to $|P_t|$, so Lagrange gives $H\cap P_t=1$.
The internal direct-product theorem
\ref{K07.04.internal_direct_product} gives
$HP_t\cong H\times P_t$ by multiplication, and $HP_t$ is normal
because both factors are. This completes the induction. At $t=s$ its
order is the product of all prime-power parts of $|G|$, hence $|G|$;
therefore the subgroup is all of $G$. This proves the displayed
multiplication isomorphism, not just an abstract cardinality equality.
Finally each $P_i$ is nilpotent by the prime-power theorem, and their
finite product is nilpotent by Proposition~\ref{prop:nilpotence-closure}.
For $G=1$ there are no prime factors: its empty product is $1$, and
the normalizer condition has no proper subgroup to test.
\end{proof}

\KEntry{K07.18.frattini_argument}{Frattini's argument}
\begin{proposition}\label{prop:frattini-argument}
Let $G$ be finite, $A\trianglelefteq G$, and $P$ a Sylow
$p$-subgroup of $A$. Then
\[
 G=A N_G(P),\qquad [G:A]=[N_G(P):N_G(P)\cap A].
\]
In particular $[G:A]$ divides $|N_G(P)|$.
\end{proposition}
\begin{proof}
For $g\in G$, normality of $A$ puts $gPg^{-1}$ in $A$, with the
same Sylow order as $P$. Sylow conjugacy inside $A$ gives $a\in A$
such that $a(gPg^{-1})a^{-1}=P$. Then $ag\in N_G(P)$ and
$g\in A N_G(P)$. The reverse inclusion is automatic. The second
group isomorphism theorem identifies
$N_G(P)/(N_G(P)\cap A)$ with $G/A$, proving the index identity
and divisibility. If $P=1$ the same argument yields $N_G(P)=G$;
thus primes not dividing $|A|$ cause no exceptional formula.
\end{proof}

\KEntry{K07.18.maximal_subgroups}{Maximal subgroups and nilpotence}
\begin{theorem}\label{thm:maximal-normal-nilpotent}
A finite group is nilpotent if and only if every maximal proper subgroup
is normal. In a nontrivial finite $p$-group every maximal proper
subgroup is normal and has index $p$.
\end{theorem}
\begin{proof}
In a nilpotent group a maximal proper subgroup $M$ satisfies
$M<N_G(M)\leq G$, so maximality gives $N_G(M)=G$, exactly
normality. Conversely suppose every maximal proper subgroup of finite
$G$ is normal. If a Sylow $P$ were not normal, $N_G(P)<G$ would
be contained in a maximal proper subgroup $M$, chosen from the finite
subgroup set. Here $M\trianglelefteq G$ and $P$ is a Sylow
subgroup of $M$. Frattini's argument gives
$G=M N_G(P)=M$, a contradiction. Thus all Sylows are normal and
the finite criterion proves nilpotence. For $G=1$ the maximal-subgroup
condition is vacuous and the conclusion still holds.

For a finite nontrivial $p$-group the prime-power theorem gives
nilpotence and hence normality of $M$. The quotient $G/M$ is a
nontrivial $p$-group with no proper nontrivial subgroup by subgroup
correspondence. Cauchy's theorem supplies an order-$p$ subgroup, which
must therefore be the whole quotient. Thus $[G:M]=p$.
\end{proof}

\KEntry{K07.18.derived_comparison}{Derived and lower central series differ}
Let $G^{(r)}$ be the derived series already defined in
\ref{K07.15.solvable_groups}. Every term is characteristic in the whole
$G$: inductively an automorphism preserving $G^{(r)}$ preserves its
self-commutator subgroup by the mixed-commutator identity. For every
$r\geq0$,
\[
 G^{(r)}\leq\gamma_{r+1}(G).
\]
Start with equality at $r=0$; if the inclusion holds, then
$[G^{(r)},G^{(r)}]\leq[\gamma_{r+1},G]=\gamma_{r+2}$.
Therefore a group of nilpotence class at most $c$ is solvable of derived
length at most $c$, by the already proved derived-series criterion.
This is a comparison theorem; it does not identify the two recurrences.
The earlier subgroup/quotient/extension tests for solvability and the
Jordan--H\"older theorem are not reproved here.

\KEntry{K07.18.nilpotence_examples}{Examples and false central-series shortcuts}
The group $Q_8$ is nonabelian and has center $\{1,-1\}$: each of
$\pm i,\pm j,\pm k$ fails to commute with one of the other two
quaternion units, using the actual matrix model in
\ref{K07.01.dihedral_quaternion_models}. Its quotient by that center
has order four and is abelian. Thus $Z_2(Q_8)=Q_8$ and its exact
class is two. Nilpotence does not imply commutativity.

In $S_3$ the commutator subgroup is $A_3$. Sign kills commutators, and
$[a,s]=a^2\ne1$ for $a=(123),s=(12)$, since $sas^{-1}=a^{-1}$.
Hence the nontrivial subgroup generated is all $A_3$.
Since $A_3$ is abelian, the second derived subgroup is trivial. But
$[A_3,S_3]=A_3$, by the same commutator and normality of $A_3$.
Consequently $\gamma_i(S_3)=A_3$ for every $i\geq2$, so $S_3$
is solvable and not nilpotent.

For a counterexample to a reverse central-series interlacing, take
$D=\langle r,s\rangle$ to be the actual symmetry group of a regular
octagon, so $|D|=16$, $r^8=s^2=1$ and $srs=r^{-1}$.
The quotient by $\langle r^2\rangle$ is abelian, while
$[r,s]=r^2$; therefore $\gamma_2(D)=\langle r^2\rangle$.
Every element is $r^b$ or $r^bs$, so it either centralizes or inverts
$r^{2a}$. Thus $\gamma_3(D)=\langle r^4\rangle$, with equality
because $[r^2,s]=r^4$, and $\gamma_4(D)=1$ because $r^4$ is
central. Its class is three. In $G=D\times C_2$, also of class three,
the extra factor $1\times C_2$ lies in $Z_1(G)$ but not in
$\gamma_2(G)=\langle r^2\rangle\times1$.
Thus $Z_1(G)\nsubseteq\gamma_2(G)$. This is why only the valid
one-way comparison in Theorem~\ref{thm:central-series-criteria} is used;
one cannot insert reverse inclusions between arbitrary upper and lower
terms.
